\section{Обзор требований к безопасности}

Данная спецификация описывает меры обеспечения численной устойчивости, потокобезопасности и защиты от деградации системы. Учитывая, что DREAM-Net допускает изменение параметров во время инференса, критически важно предотвратить нестабильность, вызванную неограниченным ростом весов, гонками данных между потоками или численными артефактами.

\textbf{Цели раздела:}
\begin{enumerate}
    \item Гарантировать ограниченность всех динамических параметров.
    \item Обеспечить корректную работу в многопоточной среде (инференс + сон).
    \item Предотвратить каскадные сбои при аномальных входных данных.
\end{enumerate}

\section{Ограничение весов (Weight Clipping)}

\subsection{Быстрые веса ($U_t$)}

Быстрые веса обновляются на каждом шаге и потенциально могут расти неограниченно при длительном воздействии сигнала с высоким $S_t$. Для предотвращения этого вводится жёсткое ограничение нормы Фробениуса:

\begin{equation}
U_t^{(l)} \leftarrow U_t^{(l)} \cdot \min\left(1, \, \frac{U_{max}}{\|U_t^{(l)}\|_F + \epsilon}\right)
\label{eq:weight_clipping}
\end{equation}
где:
\begin{itemize}
    \item $\|U\|_F = \sqrt{\sum_{i,j} U_{ij}^2}$ — норма Фробениуса
    \item $U_{max} = 10.0$ (по умолчанию, подбирается экспериментально)
    \item $\epsilon = 10^{-6}$ — численная стабильность
\end{itemize}

\textbf{Частота применения:} После каждого обновления $U_t$ (каждый шаг инференса).

\textbf{Обоснование:} Ограничение нормы предотвращает «взрыв» эффективных весов $U_t V^\top$, который мог бы привести к неограниченному росту активаций $h_t$.

\subsection{Контекстный вектор ($c_t$)}

Контекстный вектор также подлежит нормировке для предотвращения накопления ошибок связывания:
\begin{equation}
c_t \leftarrow \frac{c_t}{\max(1, \|c_t\|_2 / C_{max})}
\label{eq:context_clipping}
\end{equation}
где $C_{max} = \sqrt{d_{ctx}}$ (норма единичного вектора в пространстве FHRR).

\textbf{Частота применения:} После каждого обновления контекста (раздел~\ref{eq:context_update}).

\subsection{Статистика ошибки ($\mu_t, \sigma_t$)}

Для предотвращения деления на ноль и некорректных оценок вводятся нижние границы:
\begin{equation}
\sigma_t \leftarrow \max(\sigma_t, \, \sigma_{min})
\label{eq:sigma_floor}
\end{equation}
где $\sigma_{min} = 10^{-4}$ (по умолчанию).

\textbf{Примечание:} Верхняя граница для $\mu_t$ не требуется, так как она адаптируется к масштабу входного сигнала.

\section{Плавное восстановление после Шока (Ramp-up)}

\subsection{Мотивация}

При выходе из режима \textbf{Shock} обратные связи $W_{down}$ и контекстная модуляция были отключены. Резкое восстановление их влияния может вызвать вторичный всплеск ошибки предсказания, так как контекст за время шока не обновлялся и может быть неактуален.

\subsection{Коэффициент модуляции $\alpha_t$}

Вводится скалярный коэффициент $\alpha_t \in [0, 1]$, который плавно увеличивает влияние топ-даун сигналов:
\begin{equation}
\alpha_t = \begin{cases}
0, & \text{если режим Shock} \\
\min\left(1.0, \, \alpha_{t-1} + \frac{1}{T_{ramp}}\right), & \text{иначе}
\end{cases}
\label{eq:alpha_update}
\end{equation}
где $T_{ramp} = 50 \, \text{мс} / \Delta t$ — количество шагов рампы.

\subsection{Применение к обратным связям}

Фактическая матрица обратных связей модифицируется следующим образом:
\begin{equation}
W_{down, t}^{(l)} = \alpha_t \cdot W_{down}^{(l)}, \quad \text{для } l > 1
\label{eq:down_alpha}
\end{equation}

Аналогично модулируется вклад контекста:
\begin{equation}
\text{Modulate}_t(c_t, h_t^{(l)}) = \alpha_t \cdot \text{Modulate}(c_t, h_t^{(l)})
\label{eq:modulate_alpha}
\end{equation}

\textbf{Гарантия:} $\alpha_t$ никогда не превышает 1.0 и не опускается ниже 0.0.

\section{Потокобезопасность (Concurrency Model)}

\subsection{Архитектура двойной буферизации}

Как описано в разделе~\ref{sec:memory}, система использует два буфера весов для разделения потоков инференса и сна.

\textbf{Требования к реализации:}

\begin{table}[h]
\centering
\caption{Механизмы атомарности по платформам}
\label{tab:concurrency}
\begin{tabular}{lp{5cm}p{4cm}}
\toprule
\textbf{Платформа} & \textbf{Механизм атомарности} & \textbf{Примечание} \\
\midrule
\textbf{Python (GIL)} & Присваивание указателей атомарно & Достаточно простого swap \\
\textbf{C++} & \texttt{std::atomic<WeightBuffer*>} & Требуется memory\_order\_acquire/release \\
\textbf{Rust} & \texttt{Arc<Mutex<WeightBuffer>>} & Гарантии безопасности владения \\
\textbf{CUDA} & Копирование через отдельный stream & Избегать синхронизации с основным потоком \\
\bottomrule
\end{tabular}
\end{table}

\subsection{Инвариант целостности буферов}

Система должна гарантировать, что в любой момент времени:
\begin{equation}
\text{buffer\_active} \neq \text{buffer\_shadow}
\label{eq:buffer_invariant}
\end{equation}

\textbf{Проверка:} В отладочном режиме рекомендуется добавлять assertion перед свопом:

\begin{lstlisting}[language=Python, caption={Проверка целостности буферов}, label={lst:buffer_check}]
assert self.buffer_active is not self.buffer_shadow, "Buffer aliasing detected"
\end{lstlisting}

\subsection{Блокировка во время свопа}

Хотя своп указателей атомарен, необходимо убедиться, что ни один поток не держит ссылку на старый буфер после свопа.

\textbf{Рекомендуемая практика:}
\begin{itemize}
    \item Инференс должен запрашивать буфер в начале каждого шага \texttt{step()}.
    \item Сон должен завершить все записи до вызова \texttt{swap\_buffers()}.
    \item После свопа старый буфер (теперь shadow) может быть безопасно перезаписан.
\end{itemize}

\section{Численная устойчивость (Numerical Stability)}

\subsection{Защита от деления на ноль}

Все операции деления должны включать эпсилон-слагаемое:
\begin{equation}
\frac{a}{b} \rightarrow \frac{a}{b + \epsilon}, \quad \epsilon = 10^{-6}
\label{eq:epsilon_div}
\end{equation}

\textbf{Применяется к:}
\begin{itemize}
    \item Нормировке удивления ($\sigma_t$ в знаменателе)
    \item Нормировке весов ($\|U\|_F$ в знаменателе)
    \item Вычислению градиентов (если реализуется кастомный backward)
\end{itemize}

\subsection{Стабильность FFT/IFFT}

Операции Fourier HRR требуют внимания к численной точности:

\begin{table}[h]
\centering
\caption{Проблемы и решения для FFT}
\label{tab:fft_stability}
\begin{tabular}{lp{8cm}}
\toprule
\textbf{Проблема} & \textbf{Решение} \\
\midrule
Потеря точности в комплексной области & Использовать \texttt{torch.complex64} минимум \\
Артефакты на границах & Применять windowing (Hann/Hamming) перед FFT \\
Накопление ошибок округления & Периодическая renormalization контекста \\
\bottomrule
\end{tabular}
\end{table}

\subsection{Градиентный поток}

Все операции в DREAM-Net должны оставаться дифференцируемыми для обучения во время сна:

\begin{table}[h]
\centering
\caption{Дифференцируемость операций}
\label{tab:gradient_flow}
\begin{tabular}{lp{3cm}p{6cm}}
\toprule
\textbf{Операция} & \textbf{Дифференцируемость} & \textbf{Примечание} \\
\midrule
FFT/IFFT & ✅ Да & Поддерживается в PyTorch/JAX \\
Clipping нормы & ✅ Да & Градиент обнуляется за пределами границы \\
Indicator function ($\mathbb{I}$) & ⚠️ Нет & Использовать Straight-Through Estimator или мягкую аппроксимацию \\
Atomic swap & ✅ Н/Д & Не влияет на градиенты (операция указателей) \\
\bottomrule
\end{tabular}
\end{table}

\textbf{Рекомендация:} Для индикаторных функций в Surprise Gate использовать мягкую сигмоиду вместо жёсткого порога, чтобы сохранить градиентный поток.

\section{Защита от аномальных входов (Input Sanitization)}

\subsection{Валидация входного сигнала}

Перед обработкой каждый вход $x_t$ должен проходить базовую проверку:
\begin{equation}
\text{valid}_t = \mathbb{I}\left(\|x_t\| < X_{max}\right) \land \mathbb{I}\left(\text{isfinite}(x_t)\right)
\label{eq:input_valid}
\end{equation}
где $X_{max} = 10.0 \cdot \mu_{input}$ (адаптивный порог на основе статистики входа).

\textbf{Действия при невалидном входе:}
\begin{enumerate}
    \item Заменить $x_t$ на последнее валидное значение (hold-last).
    \item Зарегистрировать событие в логе аномалий.
    \item Не обновлять статистику $\mu_t, \sigma_t$ для этого шага.
\end{enumerate}

\subsection{Детекция NaN/Inf}

На каждом шаге рекомендуется проверять ключевые переменные:

\begin{lstlisting}[language=Python, caption={Проверка численной стабильности}, label={lst:nan_check}]
def check_numerical_stability(h, U, e, S):
    assert torch.isfinite(h).all(), "NaN in hidden state"
    assert torch.isfinite(U).all(), "NaN in fast weights"
    assert torch.isfinite(e).all(), "NaN in prediction error"
    assert 0 <= S <= 1, "Surprise out of bounds"
\end{lstlisting}

\textbf{Примечание:} В продакшн-режиме проверки могут быть отключены для производительности, но должны оставаться в отладочных сборках.

\section{Таймауты и аварийное восстановление}

\subsection{Максимальная длительность Шока}

Для предотвращения ситуации, когда система застревает в режиме \textbf{Shock} бесконечно, вводится максимальная длительность:
\begin{equation}
\text{if } (t - t_{shock\_enter}) > T_{shock\_max} \text{ then force exit to Learning}
\label{eq:shock_timeout}
\end{equation}
где $T_{shock\_max} = 500 \, \text{мс}$ (по умолчанию).

\textbf{Обоснование:} Если система не может выйти из шока естественным путём за 500 мс, вероятно, произошла деградация статистики $\mu_t, \sigma_t$. Принудительный выход позволяет сбросить состояние.

\subsection{Сброс статистики при деградации}

При детекции аномалии в статистике (например, $\sigma_t = 0$ или $\sigma_t > 1000 \cdot \sigma_{baseline}$):
\begin{enumerate}
    \item Сбросить $\mu_t, \sigma_t$ к инициализационным значениям.
    \item Сбросить быстрые веса $U_t \leftarrow U_{target}$.
    \item Зарегистрировать событие критической ошибки.
\end{enumerate}

\section{Сводная таблица инвариантов}

\begin{table}[h]
\centering
\caption{Инварианты безопасности}
\label{tab:safety_invariants}
\begin{tabular}{lp{5cm}lp{2.5cm}}
\toprule
\textbf{Инвариант} & \textbf{Формула} & \textbf{Частота проверки} & \textbf{Критичность} \\
\midrule
Норма быстрых весов & $\|U\|_F \leq U_{max}$ & Каждый шаг & Высокая \\
Норма контекста & $\|c_t\|_2 \leq C_{max}$ & Каждый шаг & Средняя \\
Коэффициент рампы & $0 \leq \alpha_t \leq 1$ & Каждый шаг & Высокая \\
Статистика ошибки & $\sigma_t \geq \sigma_{min}$ & Каждый шаг & Высокая \\
Целостность буферов & \texttt{active ≠ shadow} & При свопе & Критическая \\
Конечность значений & \texttt{isfinite(all\_tensors)} & Каждый шаг (debug) & Критическая \\
Длительность шока & $t_{shock} \leq T_{shock\_max}$ & Каждый шаг & Средняя \\
\bottomrule
\end{tabular}
\end{table}

\section{Ограничения и открытые вопросы}

Следующие аспекты безопасности требуют дальнейшей проработки:

\begin{enumerate}
    \item \textbf{Оптимальные значения $U_{max}$:} Слишком жёсткое ограничение может снизить адаптивность, слишком мягкое — допустить нестабильность.
    
    \item \textbf{Производительность проверок:} Проверка \texttt{isfinite()} на каждом шаге может добавить заметные накладные расходы на Edge-устройствах.
    
    \item \textbf{Восстановление после сброса:} Стратегия восстановления после принудительного сброса статистики не валидирована экспериментально.
    
    \item \textbf{Блокировки в многопоточности:} В реальных RTOS-системах могут потребоваться приоритеты потоков для гарантии своевременности свопа буферов.
\end{enumerate}

Эти вопросы рекомендуется включить в план тестирования на устойчивость (stress testing).

\section{Рекомендации по реализации}

\begin{enumerate}
    \item \textbf{Вынести все константы безопасности в конфигурационный файл} — это позволит настраивать пороги без перекомпиляции.
    
    \item \textbf{Реализовать режим «Safe Mode»} — при включении все проверки активны, производительность снижена, но система максимально устойчива.
    
    \item \textbf{Логировать все срабатывания защит} — количество клиппингов весов, сбросов статистики, принудительных выходов из шока. Это поможет отладке.
    
    \item \textbf{Документировать все инварианты в коде} — использовать assert с понятными сообщениями об ошибках.
\end{enumerate}
